\section{Experimental Results}
\label{sec:results}

\subsection{End-to-End Performance under Active Stress}

Figure~\ref{fig:end-to-end-frontier} and Table~\ref{tab:main-results} summarize interactive AgentDojo runs in which the harness is inserted at tool-output and tool-call boundaries. We include only batches where the undefended agent is genuinely under attack---specifically, where its security falls below 90\% and the attack succeeds at least 10\% of the time (Section~\ref{sec:stress-gate}). This filter avoids crediting a defense for episodes the base model already resisted on its own. The comparison covers 18 such batches on DeepSeek-Chat and 16 on Llama-3-70B.

In Figure~\ref{fig:end-to-end-frontier}, each point is the mean utility and security for one defense configuration. Security is the fraction of episodes in which the injected malicious goal is blocked; utility is the fraction in which the user's legitimate task is completed. Desirable defenses move toward the upper-right corner, improving safety without simply stopping all task progress.

\begin{figure*}[t]
\centering
\includegraphics[width=0.92\textwidth]{fig_end_to_end_frontier.pdf}
\caption{Utility and security on AgentDojo batches where the undefended agent is under active attack. Each point is the mean over batches for one defense. The full harness occupies the upper-right region for both model backends.}
\label{fig:end-to-end-frontier}
\end{figure*}

\begin{table*}[t]
\centering
\small
\setlength{\tabcolsep}{7pt}
\caption{Aggregate performance (\%) on AgentDojo batches where the undefended agent is under active attack. ASR is attack success rate (the complement of security). Retention is harness utility divided by no-defense utility on the same batches.}
\label{tab:main-results}
\begin{tabular}{llrrrrr}
\toprule
Model & Defense & Utility & Security & ASR & Retention & Batches \\
\midrule
DeepSeek-Chat (V3) & No defense & 67.6 & 38.0 & 62.0 & 1.00 & 18 \\
 & Spotlighting & 73.1 & 53.7 & 46.3 & 1.08 & 18 \\
 & Learned detector & 31.5 & 75.9 & 24.1 & 0.47 & 18 \\
 & \textbf{Ours} & \textbf{77.8} & \textbf{99.1} & \textbf{0.9} & \textbf{1.15} & \textbf{18} \\
\midrule
Llama-3-70B & No defense & 47.9 & 39.6 & 60.4 & 1.00 & 16 \\
 & Spotlighting & 52.1 & 53.1 & 46.9 & 1.09 & 16 \\
 & Learned detector & 52.1 & 70.8 & 29.2 & 1.09 & 16 \\
 & \textbf{Ours} & \textbf{52.1} & \textbf{100.0} & \textbf{0.0} & \textbf{1.09} & \textbf{16} \\
\bottomrule
\end{tabular}
\end{table*}

On DeepSeek-Chat, the full harness reaches 99.1\% security and 77.8\% utility, compared with 53.7\% security and 73.1\% utility for Spotlighting. Utility retention is 1.15: the harness completes more legitimate tasks than the undefended agent on the same difficult batches, while nearly eliminating successful attacks. The learned detector raises security to 75.9\%, but its 31.5\% utility and 0.47 retention show that blocking inputs aggressively can destroy task completion.

On Llama-3-70B, the full harness reaches 100.0\% security with 52.1\% utility---the same task-completion rate as Spotlighting and the learned detector, but with strictly higher security. Retention is 1.09 relative to the undefended baseline (47.9\% utility). The absolute utility level therefore reflects the open-weight model's capability on these tasks rather than excess withholding by the harness: the learned detector, which is known to over-block on other backends, does not reduce Llama utility further either. The main cross-model takeaway is that strong security transfers; utility gains depend on how much headroom the underlying agent provides.

\subsection{Performance across Attack Types}

Figure~\ref{fig:attack-type-tradeoff} breaks the full-harness results down by how the adversary frames the injection, rather than averaging over all batches. We consider four families: \emph{direct} injections embedded in external content, \emph{instruction overrides} that replace the user's goal, \emph{education-framed} attacks that disguise malicious steps as teaching material, and \emph{persona-framed} attacks that rewrite the agent's role. Batch counts per category are $(2,5,5,6)$ on DeepSeek-Chat and $(3,4,4,5)$ on Llama-3-70B.

On DeepSeek-Chat, instruction-override, education-framed, and persona-framed attacks all reach 100\% security with 80.0--83.3\% utility. Direct injections are the harder case, reaching 91.7\% security and 50.0\% utility---the adversary's instructions appear inside routine tool output without an explicit override frame, so the agent must rely on boundary checks rather than obvious goal conflict. On Llama-3-70B, security remains 100\% for all four families, while utility varies from 41.7\% for instruction overrides to 75.0\% for education-framed attacks, reflecting how each framing interacts with the model's planning behavior.

\begin{figure*}[t]
\centering
\includegraphics[width=0.92\textwidth]{fig_attack_type_tradeoff.pdf}
\caption{Utility and security of the full harness by attack framing on batches where the undefended agent is under active attack. Each point is the mean within one attack category; the dotted line marks 100\% security. Points with identical means are offset horizontally by less than one percentage point for visibility.}
\label{fig:attack-type-tradeoff}
\end{figure*}

\subsection{Utility Recovery through Boundary-Aware Repair}

The interactive sweeps above measure full task completion in live AgentDojo environments. To isolate \emph{why} the harness preserves utility, we run a separate controlled replay experiment. A single representative injected tool output is fixed, and we test a small matrix of plausible next actions---reading further, summarizing, or calling a tool---under progressively richer repair rules. This probe measures whether benign continuations can proceed after an injection is detected and stripped; it does not score end-to-end sandbox success.

Figure~\ref{fig:utility-recovery} shows how utility rises as repair components are added in sequence. A policy that withholds all suspicious content keeps 100\% security but only 20\% utility on this probe. Allowing metadata to be repaired after stripping raises utility to 60\%. Permitting benign read-like continuations after a warning raises it to 80\%. The complete communication protocol resolves the remaining failure mode and reaches 100\% utility while security stays at 100\% throughout.

\begin{figure*}[t]
\centering
\includegraphics[width=0.92\textwidth]{fig_staged_utility_recovery.pdf}
\caption{Staged utility recovery on the controlled replay probe. The curve reports utility; security remains 100\% at every stage. Plateaus between 60\%, 80\%, and 100\% indicate that some components matter only in combination rather than in isolation.}
\label{fig:utility-recovery}
\end{figure*}

Table~\ref{tab:module-ablation} complements the staged curve by removing one component at a time from the full protocol. Removing metadata repair causes the largest loss ($-40$ percentage points of utility): injected observations stay withheld instead of being cleaned and passed forward. Removing goal-aligned continuation costs 20 points, because benign actions that should follow repair remain blocked. Other single removals show no measurable change on this one scenario. A zero entry here does not mean a component is useless---the staged progression shows interaction effects that a single replay cannot separate.

\begin{table*}[t]
\centering
\small
\caption{Removing one component at a time on the controlled replay probe. Utility differences are measured relative to the full protocol; security is reported in percent.}
\label{tab:module-ablation}
\begin{tabular}{lp{0.47\columnwidth}rr}
\toprule
Removed component & Observed consequence & $\Delta U$ & Security \\
\midrule
Metadata repair & Injected observations remain withheld & $-40$ & 100 \\
Goal-aligned continuation & Benign post-repair actions remain constrained & $-20$ & 100 \\
Other single removals & No measurable change on this probe & $0$ & 100 \\
\bottomrule
\end{tabular}
\end{table*}

\subsection{Targeted Resistance to Order-Dependent Attacks}

The targeted suite evaluates a narrower question: can controlled reordering of parsed content expose attacks that look valid only in their original arrangement? Each scripted episode ($n=20$ per category) smuggles instructions across steps, fields, or observation boundaries---for example, splitting a malicious goal across reasoning stages or hiding it inside structured data. Success is measured by whether the defense blocks the attack, not by whether a full user task completes.

We compare two stages of input checking. \emph{Structural validation} applies canonical parsing and field checks only---the content must be well formed, but its order is left unchanged. \emph{Structural validation plus reordering test} additionally permutes fragments and rejects content whose meaning depends on a specific staging. Figure~\ref{fig:shuffle-gain} reports how much extra attack coverage the reordering test adds on top of structural validation alone.

Adding the reordering test improves blocked-attack rate by 100 percentage points in five categories and by 50 points in external-observation injection, where structural validation already stops half of the attacks. Together, the two stages block all 20 scripted attacks in every category.

\begin{figure*}[t]
\centering
\includegraphics[width=0.92\textwidth]{fig_shuffle_security_gain.pdf}
\caption{Additional attack coverage from adding a reordering test to structural validation on targeted order-dependent attacks ($n=20$ per category). Values are absolute percentage-point gains in the fraction of attacks blocked.}
\label{fig:shuffle-gain}
\end{figure*}

These results support a mechanism-level claim rather than a guarantee of universal robustness. They show that many evasive constructions fail once content is reorganized and rechecked, consistent with the reordering-based sensing in the Method section. Because the suite is scripted and covers a finite set of attack patterns, it validates behavior on the tested constructions rather than on arbitrary paraphrase or novel environments.
